% Macros of the blueprint. The author's own macros are at the start of each document; do not duplicate them here.

% The macros of leanblueprint, as no-ops under pdflatex, so that its dependency graph can be
% generated later: \lean{Decl.name}, \uses{label, label}, \leanok.
\newcommand{\lean}[1]{}
\newcommand{\uses}[1]{}
\newcommand{\leanok}{}

% Results that carry the paper's own number, for instance
%   \begin{cgss}{Lemma}{2.8}\label{lem:2.8} ... \end{cgss}
% and results of ours, numbered within the section, through the usual counters.
% Statements are set upright, since pf2's keywords (Assume, Prove) are small capitals, which have
% no italic shape in Computer Modern.
\theoremstyle{definition}
\makeatletter
\newtheorem*{cgss@env}{\cgss@name}
\newcommand{\cgss@name}{}
% The environment sets an anchor and makes the paper's number the text of \ref, so that
% \ref{lem:2.9} prints 2.9 and links to the lemma.
\newenvironment{cgss}[2]{\renewcommand{\cgss@name}{#1 #2}\begin{cgss@env}%
  \phantomsection\def\@currentlabel{#2}}{\end{cgss@env}}
\makeatother
\newtheorem{lemma}{Lemma}[section]
\newtheorem{proposition}[lemma]{Proposition}
\newtheorem{corollary}[lemma]{Corollary}
\newtheorem{definition}[lemma]{Definition}
\newtheorem*{hypothesis}{Hypothesis}
\theoremstyle{remark}
\newtheorem*{remark}{Remark}

% Citing assumptions as in [L95]: \asm{1} gives <1>, the assumption of the enclosing level-1
% step, and \asm[2]{0} gives <0>:2, the second assumption of the theorem itself. Steps are
% cited with pf2's \stepref{label}.
\newcommand{\asm}[2][]{$\langle#2\rangle$\if\relax\detokenize{#1}\relax\else{:}#1\fi}

% Tags of the structured proofs: a step that departs from the paper, and a step that fills a gap.
\newcommand{\new}{\textbf{[new]}}
\newcommand{\repair}{\textbf{[repair]}}

% The paper's sets of Hypothesis 2.4.
\newcommand{\Ga}{\Gamma_a}
\newcommand{\Gb}{\Gamma_b}
\newcommand{\Gc}{\Gamma_c}

% Vectors that the author's macros do not provide, in its pattern: the vectors a, b, c of Hypothesis 2.4,
% a unit vector e, and roots r, s, t.
\newcommand\Za{{\mathbf a}}
\newcommand\Zb{{\mathbf b}}
\newcommand\Zc{{\mathbf c}}
\newcommand\Ze{{\mathbf e}}
\newcommand\Zr{{\mathbf r}}
\newcommand\Zs{{\mathbf s}}
\newcommand\Zt{{\mathbf t}}

% Additional vector names for the arbitrary adjacent pair in Lemma 2.10.
\newcommand\Zp{{\mathbf p}}
\newcommand\Zq{{\mathbf q}}

% Upright Dynkin type letters, following cgss/notation.md: \typ{A}_\ell, \typ{E}_8.
\newcommand{\typ}[1]{\mathrm{#1}}

% Auxiliary results have descriptive headings, not numbers that could clash with CGSS.
% For example, \begin{namedresult}{Corollary}{of Lemma 2.8} inherits cgss's unique anchor;
% \ref prints the descriptive second argument. Existing result environments are unchanged.
\newenvironment{namedresult}[2]{\begin{cgss}{#1}{#2}}{\end{cgss}}

% The all-positive half root used in the Section 3 spanning and inclusion proofs.
\newcommand\Zh{{\mathbf h}}

% A visible citation with its exact Lean name as metadata and a pdfLaTeX-safe display.
% Unlike \path, the display may render Unicode subscripts and Greek letters as TeX math.
% Example: \leancite{LineSystem.E₈}{\texttt{LineSystem.E\ensuremath{{}_8}}}.
\newcommand{\leancite}[2]{#2}

% Local layout for the merged structured proofs; existing section settings are restored.
\newenvironment{mergedproofsection}{\begingroup\emergencystretch=3em}{\par\endgroup}

% A descriptive library-fact label (the paragraph already supplies its hyperlink anchor).
\makeatletter
\newcommand{\factlabel}[2]{\def\@currentlabel{#1}\label{#2}}
\makeatother
